\noindent
\begin{tikzpicture}[remember picture,overlay,x=1mm,y=1mm]
\begin{scope}[shift={(current page.south west)}]

% Fixed A3 canvas
\path[use as bounding box] (0,0) rectangle (420,297);

% Background
\fill[Sky] (0,0) rectangle (420,297);

% Header bar
\fill[ForestDark] (0,297) rectangle (420,202);

% Left logo
\node[anchor=west, inner sep=0pt] at (12,253)
  {\AssetImage{natureprotector-logo.png}{height=66mm,keepaspectratio}{56mm}};

% Right flag
\node[anchor=east, inner sep=0pt] at (407,258)
  {\AssetImage{flag-portugal.png}{width=38mm,keepaspectratio}{28mm}};

% Header text block
\node[anchor=north west, inner sep=0pt, align=left, text width=292mm] at (80,279)
{%
  {\MainTitle\color{Fire}NatureProtector}\\[4.2mm]
  {\SubTitle\color{white}Pipeline técnica para apoio à prevenção de incêndios}\\[5.0mm]
  {\HeaderText\color{white}\textbf{Projeto e Seminário} \;·\; LEIC \;·\; ISEL \;·\; 2025/2026}\\[3.0mm]
  {\HeaderSmall\color{white}Autores: Miguel Alves e Gabriel Mano \hspace{6mm} | \hspace{6mm} Orientadores: Artur Ferreira e Nuno Leite}\\[3.0mm]
  {\HeaderSmall\color{Fire}\textbf{A proteger florestas para gerações futuras} \hspace{3mm} · \hspace{3mm} Ideia reconhecida com o \textbf{3.º lugar} num concurso de startups}
};

% =========================================================
% BODY LAYOUT
% Menos espaço vazio no topo e no fundo
% Imagens da esquerda maiores
% Diagrama central descido para não tocar no título
% =========================================================

% Left column
\CardAt{6}{198}{116}{50}{O problema}{%
A prevenção de incêndios precisa de mais do que um valor de risco. É necessário perceber \Key{que dados chegaram}, \Key{quais foram aceites} e \Key{que evidência sustenta o estado apresentado}.%
}

% Burned area
\draw[fill=white, draw=Line, line width=0.35mm, rounded corners=2.4mm]
  (6,142) rectangle ++(116,-60);
\node[anchor=center, inner sep=0pt] at (64,112)
  {\AssetImage{burned-area.png}{width=110mm,height=54mm,keepaspectratio}{48mm}};

% Imagem de risco
\draw[fill=white, draw=Line, line width=0.35mm, rounded corners=2.4mm]
  (6,76) rectangle ++(116,-69);
\node[anchor=center, inner sep=0pt] at (64,41.5)
  {\AssetImage{perigo-incendio.jpeg}{width=110mm,height=60mm,keepaspectratio}{54mm}};

% Centre column
\draw[fill=white, draw=Line, line width=0.35mm, rounded corners=2.4mm]
  (126,198) rectangle ++(202,-120);

% Título dentro da caixa
\node[anchor=north west, inner sep=0pt] at (130,194)
  {{\SectionText\color{Navy}Como funciona}};

% Diagrama central mais baixo para não tapar o título
\node[anchor=center, inner sep=0pt] at (227,139)
  {\AssetImage{architecture-overview.png}{width=188mm,height=86mm,keepaspectratio}{78mm}};

% Flow row compacta e contida dentro da caixa "Como funciona"
\FlowChipAt{146}{91}{Forest}{sensores}
\node[color=ForestDark] at (162,91) {\small$\rightarrow$};

\FlowChipAt{178}{91}{Navy}{eventos}
\node[color=ForestDark] at (194,91) {\small$\rightarrow$};

\FlowChipAt{219}{91}{FireDark}{processamento}
\node[color=ForestDark] at (246,91) {\small$\rightarrow$};

\FlowChipAt{268}{91}{ForestDark}{projeções}
\node[color=ForestDark] at (290,91) {\small$\rightarrow$};

\FlowChipAt{310}{91}{Navy}{visualização}

% Ideia central mais baixa e maior
\FireCardAt{126}{72}{202}{65}{Ideia central}{%
O sistema não tenta provar uma previsão científica final. Demonstra uma \OrangeKey{pipeline rastreável}: cada leitura deixa um percurso observável desde o evento até às projeções e alertas.%
}

% Right column
\draw[fill=white, draw=Line, line width=0.35mm, rounded corners=2.4mm]
  (332,198) rectangle ++(82,-121);

% título dentro da caixa
\node[anchor=north west, inner sep=0pt] at (336,194)
  {{\SectionText\color{Navy}Tecnologias}};

% grelha 2x4
\TechAt{352}{178}{dotnet.png}{.NET}
\TechAt{392}{178}{csharp.png}{C\#}

\TechAt{352}{154}{rabbitmq.png}{RabbitMQ}
\TechAt{392}{154}{postgreSQL.png}{PostgreSQL}

\TechAt{352}{130}{influxdb.png}{InfluxDB}
\TechAt{392}{130}{grafana.jpg}{Grafana}

\TechAt{352}{106}{react-tsx.png}{React/TSX}
\TechAt{392}{106}{docker.png}{Docker}

% contributo
\CardAt{332}{71}{82}{64}{Contributo}{%
\begin{itemize}
  \item eventos simulados;
  \item processamento durável;
  \item rejeições e quarentenas;
  \item avaliações candidatas de risco;
  \item API e observabilidade.
\end{itemize}%
}

\end{scope}
\end{tikzpicture}

% exactly one output page
\mbox{}